\chapter{热力学定律}
\setcounter{example_counter}{1}

\section{热力学第一定律}

\subsection{基本概念}

\subsubsection{过程与途径}

系统从一个状态变化到另一个状态，称为一个\textbf{过程}，完成这个过程的具体方式称为\textbf{途径}。

若系统的变化非常缓慢，任意状态都可近似看作平衡态，则称为\textbf{准静态过程}。

\subsubsection{状态量与过程量}

\begin{itemize}
    \item 状态量与具体途径无关，$p,V,T,E$都是状态量，其变化只与系统的始末状态有关
    \item 过程量与具体途径有关，$A,Q$是过程量，用不同的途径完成同一个过程，做功和吸热可能不同
\end{itemize}

\subsection{热力学第一定律}

\subsubsection{表达式}

在一给定的过程中，外界对系统做的功和传给系统的热量之和，等于系统内能的增量。

\begin{formula}
    $Q=\Delta E+A$
\end{formula}

\begin{itemize}
    \item $Q$表示系统所吸收的热量，规定吸热为正，放热为负；
    \item $A$表示系统对外界所做的功，规定体积膨胀为正，体积收缩为负
\end{itemize}

\warn{热力学第一定律有不同的表达形式，务必注意$Q,A$符号的规定。}

\note{为了和高中写法相区分，这里做功采用多数大学教材所使用的符号$A$。}

\subsubsection{$\Delta E$的计算}

内能$E=\dfrac{i}{2}\nu RT$是状态量，$\Delta E$只与初末状态的温度变化有关，与具体途径无关。

对于一定气体而言，内能的变化$\Delta E=\dfrac{i}{2}\nu R\Delta T$

另外，结合$E=\dfrac{i}{2}pV$，可知当$p$不变时，$\Delta E=\dfrac{i}{2}p\Delta V$；当$V$不变时，$\Delta E=\dfrac{i}{2}V\Delta p$

\newpage

\subsubsection{$A$的计算}

气体的体积从$V_1$变为$V_2$的过程中，对外所做的功可通过积分求解。

\begin{formula}
    $A=\displaystyle \int_{V_1}^{V_2}{p~dV}$
\end{formula}

在$p-V$图像中，$A$的绝对值就等于曲线与横轴围成的面积，膨胀为正，收缩为负。

\begin{figure}[htb]
    \centering
    \includegraphics[width=0.35\linewidth]{asset/ch7/T1.png}
\end{figure}

\begin{example}{热力学第一定律}
    
\end{example}

\subsection{热容}

\subsubsection{热容的概念}

$\SI{1}{mol}$物质吸收的热量与升高温度之比称为\textbf{摩尔热容}。

定体摩尔热容$C_{V,m}$和定压摩尔热容$C_{p,m}$的表达式如下：

\begin{formula}
    $C_{V,m}=\dfrac{i}{2}R,~~C_{p,m}=\dfrac{i+2}{2}R,~~C_{p,m}-C_{V,m}=R$
\end{formula}

\note{定压摩尔热容较大，这是因为等压条件下，气体吸热膨胀，对外做功，升温更难。}

\vspace{0.3em}

定压摩尔热容$C_{p,m}$与定体摩尔热容$C_{V,m}$的比$\gamma=\dfrac{i+2}{i}$称为\textbf{比热比}。

\subsubsection{$Q$的计算}

在等体或等压过程中，热量可采用相应的摩尔热容$C_m$计算。

\begin{formula}
    $Q=\nu C_m \Delta T$
\end{formula}

\note{其他过程当然也可能有热量收放，但不能直接用公式计算，需要用热力学第一定律求解。}

\begin{example}{热容}
    
\end{example}

\subsection{等值过程}

\subsubsection{等体过程}

在等体过程中，体积$V$不变，因此$A=0$.~热量$Q$使用定体摩尔热容计算，$Q=\dfrac{i}{2}\nu R\Delta T$

内能变化$\Delta E=\dfrac{i}{2}\nu R\Delta T$，与热力学第一定律$Q=\Delta E$契合。

\subsubsection{等压过程}

在等压过程中，压强$p$不变，因此$A=p\Delta V$.~热量$Q$使用定压摩尔热容计算，$Q=\dfrac{i+2}{2}\nu R\Delta T$

内能变化$\Delta E=\dfrac{i}{2}\nu R\Delta T$，与热力学第一定律$Q=\Delta E+A$契合。

\subsubsection{等温过程}
在等温过程中，温度$T$不变，因此$\Delta E=0$.~根据理想气体状态方程，$pV=\nu RT$为定值。

代入做功的公式，$A=\displaystyle\int_{V_1}^{V_2}{pdV}=p_1V_1\ln{\dfrac{V_2}{V_1}}$，利用热力学第一定律可求出热量$Q=A$

\warn{等温过程虽然温度不变，但也有热量收放。}

\subsubsection{绝热过程}

在绝热过程中，没有热量收放，$Q=0$.~绝热过程的特点是$pV^{\gamma}$为定值，其中$\gamma=\dfrac{i+2}{i}$为比热容比。

\begin{formula}
    绝热过程$pV^{\gamma}=$定值
\end{formula}

代入做功的公式，$A=\displaystyle\int_{V_1}^{V_2}{pdV}=\dfrac{1}{\gamma-1}(p_1V_1-p_2V_2)$

在图像上，绝热过程比等温过程更陡峭，这是因为绝热膨胀气体对外做功，内能减小使温度下降，压强减小更多。

\subsubsection{等值过程的对比}

\subsection{气体向真空的绝热自由膨胀}

气体向真空的绝热自由膨胀不是准静态过程，不吸热、不做功，因此内能不变，温度不变。

\note{但这一过程不是准静态过程，所以不能说是等温过程。}

\section{循环过程与热机}

\subsection{循环过程的特点}

循环过程是指系统经一系列变化后又回到初态的整个过程，可用$p-V$图上的一个闭合曲线表示，循环一周内能不变。

\subsection{热循环与制冷循环}

\subsubsection{热循环}

顺时针方向称为热循环，每循环一周，系统吸热并对外做功。

循环曲线所围面积就是系统对外所做的净功$A$，设吸热$Q_1$，放热$|Q_2|$，定义热循环效率为

\begin{formula}
    $\eta=\dfrac{A}{Q_1}=1-\dfrac{|Q_2|}{Q_1}$
\end{formula}

\subsubsection{制冷循环}

逆时针方向称为制冷循环，每循环一周，外界对系统做功，系统放热。

循环曲线所围面积是外界对系统做的净功$|A|$，设吸热$Q_1$，放热$|Q_2|$，定义制冷系数为

\begin{formula}
    $w=\dfrac{Q_2}{A}$
\end{formula}

\note{热循环效率一定在$0\sim 1$之间，但制冷系数可以大于$1$.}

\subsection{卡诺热机}

卡诺循环是工质只和两个恒温热源交换热量的无摩擦的准静态循环，由等温过程和绝热过程组成。

卡诺热机是按卡诺循环工作的热机，卡诺热机的效率是热循环效率的上限，只与两个热源的温度有关。

\note{卡诺热机按逆向循环就是卡诺制冷机，其制冷系数也是上限。}

\begin{formula}
    $\eta=1-\dfrac{T_2}{T_1},~~w=\dfrac{T_2}{T_1-T_2}$
\end{formula}

\section{热力学第二定律}

\subsection{热力学第二定律的宏观表述}

热力学第二定律是关于自然过程方向性的一条基本定律，有两种等价的表述。

\begin{itemize}
    \item 开尔文表述：其\textbf{唯一效果}是热量全部转变为功的过程是不可能的。（第二类永动机不可制成）
    \item 克劳修斯表述：热量不能\textbf{自动地}从低温物体传向高温物体。

\end{itemize}

\subsection{热力学第二定律的微观含义}

\subsubsection{可逆过程}

可逆过程是结果（系统和外界的变化）可完全被消除的过程，可逆过程必然可以沿原路径的反向进行，结果是系统和外界同时复原。

卡诺循环就是可逆过程，但一切与热现象有关的实际宏观过程（自然过程）都不可逆。

\subsubsection{热力学概率}

系统的同一个宏观状态实际上可能对应于非常多的微观状态，孤立系统的各个微观状态出现的概率相等。

任一宏观态所对应的微观态数称为该宏观态的热力学概率$\Omega$，一个孤立系统其内部自发进行的过程，总是由热力学概率小的宏观态向热力学概率大的宏观态过渡。

\subsubsection{玻尔兹曼熵公式}

玻耳兹曼引入熵$S$表示平衡态系统无序性的大小，熵是系统状态的函数。

\begin{formula}
    $S=k\ln{\Omega}$
\end{formula}

\subsubsection{熵增原理}

孤立系统内的自然过程总是沿着熵增加的方向进行，这是热力学第二定律的另一种表述。

即对于孤立系统中进行的过程，$\Delta S\geq 0$，只有可逆过程的$\Delta S=0$，自然过程的$\Delta S$必然大于零。

\note{可逆绝热过程就是熵变为零的过程。}

\section{拓展内容}

\subsection{熵变}

\subsubsection{克劳修斯熵公式}

一个系统进行可逆过程时，熵变计算公式为

\begin{formula}
    $dS=\dfrac{dQ}{T},~~~\Delta S=\displaystyle\int_1^2{\dfrac{dQ}{T}}$
\end{formula}

\note{其中$Q$表示气体吸收的热，熵$S$的单位为$\unit{J/K}$.}

对于$\SI{1}{mol}$理想气体而言，利用理想气体状态方程，可得

\begin{formula}
    $\Delta S=C_{V,m}\ln{\dfrac{T_2}{T_1}}+R\ln{\dfrac{V_2}{V_1}}$
\end{formula}

\subsubsection{熵变的计算}

熵是状态函数，熵变只与始末状态有关，而与变化过程无关。

因此要求熵变，则不用考虑实际的变化过程，只需设计一个可逆过程，利用克劳修斯熵公式计算即可。

